\section{Memory、continual learning、interpretability 与 alignment}

上一章讨论模型对当前世界的错误表示，本章追问系统如何跨时间保存信息、更新参数，并在新环境中保持目标一致。Memory、continual learning、interpretability 与 alignment 经常一起出现，因为它们共同处理“系统长期运行后会变成什么”。

\subsection{Long-term memory：不只是把上下文拉长}

长期 memory 至少包含写入、检索、更新、冲突处理、压缩、遗忘与权限。可以把 memory store 记为 $M_t$，新 observation 为 $o_t$，更新策略为
\[
M_{t+1}=U(M_t,o_t;\rho,\kappa),
\]
$\rho$ 是保留/淘汰策略，$\kappa$ 是冲突解决规则。若新事实与旧事实矛盾，系统可覆盖、保留版本、降低置信度或请求确认；没有 policy 的“全部保存”会导致噪声和隐私风险。

\lecturefigure{slide-139.jpg}{Long-term memory 的目标：跨会话保存、检索与个性化}{CS25 V6 official deck, p.~139}

\lecturefigure{slide-140.jpg}{Memory 方法：向量库、摘要、压缩、更新与遗忘}{CS25 V6 official deck, p.~140}

\teachervoice{01:03:52--01:06:32，老师强调 memory 的难点不在“能否存”，而在 contradiction、pruning 与 relevance。实践经验：unbounded context 不等于可靠 memory manager。}

\begin{lstlisting}[language=Python,caption={带冲突与淘汰策略的 memory update 伪代码}]
def update_memory(memory, observation, policy):
    candidates = retrieve_conflicts(memory, observation)
    resolution = policy.resolve(observation, candidates)
    memory.apply(resolution)
    if memory.size > policy.capacity:
        memory.evict(policy.rank_for_eviction(memory))
    return memory
\end{lstlisting}

\begin{knowledgebox}{Memory 术语消化}
\textbf{Context memory} 是当前 prompt 中的 token；\textbf{episodic memory} 保存过去事件；\textbf{semantic memory} 保存事实或概念；\textbf{vector store} 是检索实现；\textbf{parameter memory} 是写入权重的统计规律。它们更新速度、来源追踪和遗忘机制不同。
\end{knowledgebox}

\subsection{Lifelong learning：什么才算模型真的学了}

课程采用较严格定义：RAG、context memory 与 distillation 可以改善持续使用，但若模型参数与技能结构没有在线适应，不一定叫 true continual learning。人类能在连续交互中更新，而 LLM 常依赖周期性离线再训练。关键指标包括 plasticity、retention、forward transfer、backward transfer 与 catastrophic forgetting。

\lecturefigure{slide-141.jpg}{Lifelong learning like humans：部署后持续学习的目标与困难}{CS25 V6 official deck, p.~141}

\lecturefigure{slide-142.jpg}{Current work vs. needed：prompt memory、distillation、model editing 与真正持续适应}{CS25 V6 official deck, p.~142}

\teachervoice{01:06:32--01:08:08，老师明确区分 memory/RAG 与 parameter-level continual learning。课堂提示：把资料放进向量库可以更新知识源，但没有改变模型本身的学习动力学。}

\subsection{Model editing：局部改权重为何仍会外溢}

Rank-One Model Editing（ROME）等方法试图通过低秩更新修改某个事实。抽象地，若某层权重为 $W$，更新为
\[
W'=W+uv^{\top},
\]
$u,v$ 控制写入方向。局部低秩不等于语义局部：同一参数参与多个上下文，修改一个关系可能影响 paraphrase、相关事实或无关行为。评估需要 efficacy、generalization、specificity 与 portability，而不是只看目标 prompt。

\lecturefigure{slide-143.jpg}{Model editing for continual learning：局部事实写入及其副作用}{CS25 V6 official deck, p.~143}

\begin{warningbox}{Model editing 不是数据库 UPDATE}
权重是分布式表示，没有独立“事实槽位”。即使目标问题改对，也可能破坏邻近概念、只对一种措辞有效，或与后续训练冲突。编辑系统必须记录 provenance、版本和 rollback，而不是把单次成功当作安全在线学习。
\end{warningbox}

长期系统还需要定义 memory 与 parameter update 的优先级。高时效、可撤销、带来源的信息更适合外部 store；反复使用、稳定且广泛迁移的规律才可能值得写入参数。在线更新前应在 shadow model 验证 locality、retention 与 safety，再通过版本化 rollout 发布。若 retrieval、memory 和权重各自保存不同版本的事实，系统必须规定冲突时谁拥有 authority，并向用户暴露来源。持续学习的难点因此不是一次 gradient step，而是完整 lifecycle：采集、验证、写入、监控、回滚与遗忘。

\subsection{Interpretability：理解内部机制而非只看输出}

Interpretability 试图识别模型表示了什么、哪些 circuit 促成输出、失败能否被提前检测。Behavioral eval 只能看到输入输出；mechanistic interpretability 进一步分析 activation、feature、attention/circuit 与 intervention。真正强证据来自改变内部变量后输出按预测变化，而不是仅做相关性可视化。

\lecturefigure{slide-145.jpg}{LLM interpretability：从 black box 到 feature、circuit 与 failure detection}{CS25 V6 official deck, p.~145}

\begin{importantbox}{解释的证据梯度}
可视化 attention 是描述；probe 是可解码性证据；activation patching/ablation 是干预证据；跨样本、跨模型复现才支持更强机制主张。解释方法也可能遗漏分布式、冗余或 context-dependent computation。
\end{importantbox}

\subsection{Alignment problem：目标、行为与新环境之间的缝隙}

Alignment 要求系统在不同环境中追求正确目标并遵守约束。形式上，训练 reward $R_{\mathrm{train}}$ 只是人类意图 $U$ 的 proxy。若 policy 最大化
\[
\pi^{\star}=\arg\max_{\pi}\mathbb{E}[R_{\mathrm{train}}],
\]
但 $R_{\mathrm{train}}\neq U$，就可能出现 reward hacking、shortcut、deception 或 distribution shift。问题不是模型“不听话”这么简单，而是目标表示与监督覆盖不完整。

\lecturefigure{slide-146.jpg}{The alignment problem：模型优化目标与人类价值的差距}{CS25 V6 official deck, p.~146}

\lecturefigure{slide-147.jpg}{Failure modes：reward hacking、deception、shortcut 与 distribution shift}{CS25 V6 official deck, p.~147}

\teachervoice{01:08:58--01:11:00，老师把 alignment risk 分成 shortcut、reward hacking、deception、post-hoc explanation 与新环境失配。课堂提示：训练集上“表现正确”不等于目标被正确内化。}

\subsection{Alignment techniques：不同方法约束不同 failure channel}

Human/AI feedback 调整输出偏好；process supervision 检查中间步骤；constitutional AI 用规则生成批评与修订；scalable oversight 试图让有限监督覆盖更强系统；interpretability 则寻找内部危险机制。它们互补而非替代：规则可能冲突，AI feedback 可能继承偏见，process label 可能奖励表面格式，interpretability 也可能不完整。

\lecturefigure{slide-148.jpg}{Alignment techniques 一：feedback、preference、interpretability 与 process supervision}{CS25 V6 official deck, p.~148}

\lecturefigure{slide-149.jpg}{Alignment techniques 二：constitutional AI、scalable oversight 与 human feedback}{CS25 V6 official deck, p.~149}

\teachervoice{01:11:00--01:13:02，老师把 constitutional AI、scalable oversight、process supervision 与 interpretability 放在同一张工具箱中，并明确指出 human feedback 本身并不保证 alignment。}

\begin{knowledgebox}{Alignment 方法矩阵}
\begin{tabularx}{\textwidth}{p{0.2\textwidth}p{0.25\textwidth}X}
\toprule
方法 & 主要约束对象 & 仍未解决的问题 \\
\midrule
Preference feedback & 输出排序与风格 & 偏好代表性、reward hacking \\
Process supervision & 中间步骤 & 多条等价路径、标注成本、表面迎合 \\
Constitutional AI & 显式原则 & 规则冲突、解释与执行间差距 \\
Scalable oversight & 强系统监督 & 弱监督者如何可靠判断复杂任务 \\
Interpretability & 内部 feature/circuit & 覆盖不完整、解释可误导 \\
\bottomrule
\end{tabularx}
\end{knowledgebox}

Alignment eval 也需要分层。Capability eval 测系统能否完成任务，alignment eval 测它在目标冲突、诱导、权限边界和新分布下是否仍选择允许行为；二者不可互相替代。测试集应包含正常请求、adversarial prompt、工具返回恶意内容、长轨迹中的局部诱惑以及 shutdown/abstention 场景。还要报告 false refusal，因为把所有困难请求拒绝可以提高某些 safety 指标，却损害可用性。最终需要画出 helpfulness--safety--cost 的 Pareto frontier，而不是用单一“安全分”隐藏权衡。

\subsection{本章小结}

Memory 管理外部状态，continual learning 改变长期适应，model editing 尝试局部写权重，interpretability 提供内部证据，alignment 约束目标与行为。它们不能互相替代：一个有检索库的系统未必会持续学习，一个能解释 attention 的系统也未必安全。

